\documentclass[tikz,border=10pt]{standalone}
\usetikzlibrary{arrows.meta, positioning, shapes.geometric, shapes.misc, fit, backgrounds, calc, matrix}

\tikzset{
    >=Latex,
    font=\sffamily\large,
    node distance=1.5cm and 3cm,
    dag_result/.style={
        rectangle, 
        draw=green!70!black, 
        fill=green!15, 
        very thick, 
        minimum width=4.5cm, 
        minimum height=1.5cm, 
        text width=5cm,
        align=center
    },
    dag_lemma/.style={
        rounded rectangle, 
        draw=orange!80!black, 
        fill=yellow!20, 
        very thick, 
        minimum width=4.5cm, 
        minimum height=1.5cm, 
        text width=5cm,
        align=center
    },
    dag_model/.style={
        ellipse, 
        draw=blue!70!black, 
        fill=blue!12, 
        very thick, 
        minimum width=4.5cm, 
        minimum height=1.5cm, 
        text width=4.5cm,
        align=center
    },
    dag_unformalized/.style={
        rectangle, 
        draw=gray!60!black, 
        fill=gray!10, 
        very thick, 
        dashed, 
        minimum width=4.5cm, 
        minimum height=1.5cm, 
        text width=5cm,
        align=center
    },
    dag_conditional/.style={
        rectangle, 
        rounded corners,
        draw=orange!80!black, 
        fill=orange!10, 
        very thick, 
        minimum width=4.5cm, 
        minimum height=1.5cm, 
        text width=5cm,
        align=center
    },
    dag_partial/.style={
        rectangle,
        rounded corners,
        draw=yellow!55!black,
        fill=yellow!10,
        very thick,
        dashed,
        minimum width=4.5cm,
        minimum height=1.5cm,
        text width=5cm,
        align=center
    },
    dag_scaffold/.style={
        rectangle,
        draw=gray!55!black,
        fill=white,
        very thick,
        dotted,
        minimum width=4.5cm,
        minimum height=1.5cm,
        text width=5cm,
        align=center
    },
    dag_caveat/.style={
        diamond, 
        draw=red!70!black, 
        fill=red!10, 
        minimum width=4cm, 
        align=center, 
        inner sep=0pt, 
        font=\sffamily\small
    },
    dag_arrow/.style={
        -{Stealth[scale=1.2]}, 
        very thick,
        shorten >=1.5pt,
        shorten <=1.5pt
    },
    dag_dashed_arrow/.style={
        -{Stealth[scale=1.2]}, 
        very thick, 
        dashed,
        shorten >=1.5pt,
        shorten <=1.5pt
    },
    dag_template_legend/.style={
        minimum width=2.6cm,
        minimum height=0.8cm,
        text width=2.45cm,
        inner sep=3pt,
        font=\sffamily\small,
        align=center
    },
    dag_template_note/.style={
        rectangle,
        draw=gray!50,
        fill=gray!5,
        rounded corners,
        text width=13cm,
        align=left,
        font=\sffamily\small,
        inner sep=6pt
    }
}

% Legend Helper
\newcommand{\daglegend}[2]{
    \node[draw, rounded corners, very thick, fit=#1, inner sep=12pt, label=above:\textbf{\Large Legend}] (#2) {};
}


\begin{document}
\begin{tikzpicture}[node distance=1.1cm and 2.2cm,
  every node/.style={font=\sffamily\small}]

\dagPaperMetadata{Driver Surge Pricing}{Nikhil Garg and Hamid Nazerzadeh}{Management Science / arXiv v4}{2021}{https://arxiv.org/pdf/1905.07544}
\dagPaperLegendRightOfMetadata{
  \node[dag_model, dag_template_legend] (legModel) {formalized\\model};
  \node[dag_lemma, right=0.35cm of legModel, dag_template_legend] (legLemma) {formalized\\lemma};
  \node[dag_result, right=0.35cm of legLemma, dag_template_legend] (legResult) {formalized\\result};
}
\daglegend{(legModel)(legLemma)(legResult)}{Legend}

\begin{dagPaperBody}
\begin{scope}[yshift=-0.8cm]

\node[dag_model, text width=4.8cm] (Model) at (0,0) {
  \textbf{Model and IC definitions}\\
  measurable trip-length policies; single-state and two-state dynamic incentive compatibility
};

\node[dag_lemma, text width=4.7cm] (Renewal) at (6.2,0) {
  \textbf{Section 2.2 renewal reward}\\
  lifetime hourly earnings equal expected cycle reward divided by expected cycle time
};

\node[dag_lemma, text width=4.7cm] (Lemma1) at (12.4,0) {
  \textbf{Lemma 1}\\
  dynamic earnings decompose into state reward rates weighted by long-run state time
};

\node[dag_lemma, text width=4.5cm] (Lemma2) at (18.6,0) {
  \textbf{Lemma 2}\\
  two-state CTMC switch probability has the stated closed form
};

\node[dag_lemma, text width=4.7cm] (Lemma3) at (24.8,0) {
  \textbf{Lemma 3}\\
  long-run state-time fractions equal the renewal-cycle time-fraction formula
};

\node[dag_result, text width=4.7cm] (Thm1) at (0,-4.4) {
  \textbf{Theorem 1}\\
  measurable single-state optimal policies are threshold policies
};

\node[dag_result, text width=4.7cm] (Prop31) at (6.2,-4.4) {
  \textbf{Proposition 3.1}\\
  affine single-state prices with $0 \le a \le m/\lambda$ are incentive compatible
};

\node[dag_lemma, text width=4.7cm] (Remarks) at (18.6,-4.4) {
  \textbf{Remarks 1, 3, and 4}\\
  $q(u)/u$ is decreasing, tends to the switch rate at zero, and gives nonnegative switch-time slack
};

\node[dag_result, text width=4.7cm] (Lemma4) at (0,-8.8) {
  \textbf{Lemma 4}\\
  reward-cutoff optimal threshold exists and is unique up to null off-boundary sets
};

\node[dag_lemma, text width=4.7cm] (Lemma5) at (6.2,-8.8) {
  \textbf{Lemma 5}\\
  fixed-response variational proof gives a.e. canonical policy forms
};

\node[dag_lemma, text width=4.7cm] (Lemma6) at (12.4,-8.8) {
  \textbf{Lemma 6}\\
  endpoint reward derivatives have the sign of the normalized response
};

\node[dag_lemma, text width=4.7cm] (Lemmas78) at (18.6,-8.8) {
  \textbf{Lemmas 7--8}\\
  affine positive-additive responses are quasi-convex; negative-additive responses are quasi-concave
};

\node[dag_lemma, text width=4.7cm] (Lemmas910) at (24.8,-8.8) {
  \textbf{Lemmas 9--10}\\
  structured surge and non-surge ratio intervals make endpoint derivatives positive
};

\node[dag_result, text width=4.8cm] (Thm4) at (18.6,-13.2) {
  \textbf{Theorem 4}\\
  optimal policies have the listed non-surge/surge forms, with a.e. representatives
};

\node[dag_result, text width=5.1cm] (Thm2) at (30.8,-8.8) {
  \textbf{Theorem 2}\\
  multiplicative pricing yields reject-long/reject-short optimal-policy forms; an explicit instance has positive finite cutoff deviations in both states
};

\node[dag_result, text width=5.1cm] (Thm3) at (30.8,-13.2) {
  \textbf{Theorem 3}\\
  structured CTMC prices give feasible sequential current-bounds IC and a.e. accept-all uniqueness
};

% Verified support edges.
\draw[dag_arrow] (Model.east) -- (Renewal.west);
\draw[dag_arrow] (Renewal.east) -- (Lemma1.west);
\draw[dag_arrow] (Lemma2.east) -- (Lemma3.west);
\draw[dag_arrow] (Model.south) -- (Thm1.north);
\draw[dag_arrow] (Thm1.east) -- (Prop31.west);
\draw[dag_arrow] (Thm1.south) -- (Lemma4.north);
\draw[dag_arrow] (Lemma4.east) -- (Lemma5.west);
\draw[dag_arrow] (Lemma5.east) -- (Lemma6.west);
\draw[dag_arrow] (Lemma2.south) -- (Remarks.north);
\draw[dag_arrow, rounded corners=6pt]
  (Lemma1.south) -- ++(0,-2.15) -| (Lemma6.north);
\draw[dag_arrow, rounded corners=6pt]
  (Remarks.south west) -- ++(-0.25,-1.15) -| (Lemma6.north east);
\draw[dag_arrow] (Lemma6.east) -- (Lemmas78.west);
\draw[dag_arrow] (Lemmas78.east) -- (Lemmas910.west);
\draw[dag_arrow] (Lemma3.south) -- (Lemmas910.north);
\draw[dag_arrow, rounded corners=6pt]
  (Remarks.south east) -- ++(0.65,-1.15) -| (Lemmas910.north);
\draw[dag_arrow, rounded corners=6pt]
  (Lemma5.south east) -- ++(0,-0.85) -| (Thm4.north west);
\draw[dag_arrow] (Lemmas78.south) -- (Thm4.north);
\draw[dag_arrow] (Lemmas910.south west) -- (Thm4.north east);
\draw[dag_arrow] (Lemma3.south east) to[out=-25,in=130] (Thm2.north);
\draw[dag_arrow] (Thm4.east) to[out=15,in=-145] (Thm2.south west);
\draw[dag_arrow] (Lemmas910.south east) -- (Thm3.north west);
\draw[dag_arrow, rounded corners=6pt]
  (Thm4.south) -- ++(0,-0.85) -| (Thm3.south);

\end{scope}
\end{dagPaperBody}
\end{tikzpicture}
\end{document}
